\begin{filecontents*}[overwrite]{corl_2026.sty}
% File: corl_2026.sty 
%
% Latex templates for the Conference on Robot Learning (CoRL)
%
% This template is heavily inspired by the NeurIPS, ICML, ICLR and IEEE Transactions latex templates.
% Hence we would like to thank: Roman Garnett and the previous mantainers of the NIPS style, Percy Liang and the previous mantainers of the ICML style, Hugo Larochelle for the ICLR style, and Michael Shell for the IEEE Transactions style.
%
% History:
% 2017/04/16 - First revision by Roberto Calandra (roberto.calandra@berkeley.edu).
%	       Main changes:
%	       - The abstract is more compact compared to NeurIPS/ICML
%	       - References are by default using natbib with squared numbers (e.g., [1])
%	       - DOI fields from the bibtex are automatically converted to hyperlinks 
%	       to the corresponding page
%	       - acknowledgments are now a command, and the corresponding subsubsection is
%	       automatically included only in the final version
% 2017/06/12 - Modified to use corlabbrvnat.bst, which order the reference by order of appearance in the paper
% 2017/06/13 - fixed typo
% 2018/05/09 - Slightly modified for CoRL 2018 by Jun Morimoto (xmorimo@atr.jp)
% 2019/01/28 - Slightly modified for CoRL 2019 by Jun Nakanishi (jnakanis@meijo-u.ac.jp)
% 2020/02/02 - Slightly modified for CoRL 2020 by Cynthia Matuszek (cmat@umbc.edu)
% 2020/08/19 - Added preprint option by Roberto Calandra (rcalandra@fb.com)
% 2021/05/06 - Slightly modified for CoRL 2021 by Gerhard Neumann (gerhard.neumann@kit.edu)
% 2022/03/09 - Slightly modified for CoRL 2022 by Minas Liarokapis (minas.liarokapis@auckland.ac.nz)
% 2022/03/06 - Slightly modified for CoRL 2023 by Marc Toussaint (toussaint@tu-berlin.de)
% 2024/03/26 - Slightly modified for CoRL 2024 by David Held (dheld@andrew.cmu.edu)
% 2026/01/10 - Slightly modified for CoRL 2026 by Yoonchang Sung (yoonchang.sung@ntu.edu.sg)
%
% TODO: nohyperref is not working at the moment
% 

\NeedsTeXFormat{LaTeX2e}

% Content to be changed from year to year
\ProvidesPackage{corl_2026}[2026/08/15 CORL2026 submission/preprint/camera-ready style file]
\newcommand{\@conferenceordinal}{10th}
\newcommand{\@conferenceyear}{2026}
\newcommand{\@conferencelocation}{Austin TX, USA}


%%%%%%%%%%%%%%%%%%%%%%%%%%%%%%%%%%%%%%%%%%%%%%%%%%%%%%%%%%%%%%%%%%%%%%%%
% Accepted options: [final,preprint,nonatbib,nohyperref]

% Declare the final option, which creates camera-ready copy
\newif\if@conferencefinal\@conferencefinalfalse
\DeclareOption{final}{
  \@conferencefinaltrue
}

% Declare the preprint option, which creates a camera-ready copy without the corl footnote
\newif\if@preprinttype\@preprinttypefalse
\DeclareOption{preprint}{
  \@preprinttypetrue
}

% The natbib package is loaded by default. Declaring the nonatbib option, does not load natbib in case of package clash (users can pass options to natbib via \PassOptionsToPackage)
\newif\if@natbib\@natbibtrue
\DeclareOption{nonatbib}{
  \@natbibfalse
}

% The hyperref package is loaded by default. Declaring the  nohyperref option, does not load the hyperref.
\DeclareOption{nohyperref}{%
  \gdef\nohyperref{1}
}

% Activate the options
\ProcessOptions\relax

%%%%%%%%%%%%%%%%%%%%%%%%%%%%%%%%%%%%%%%%%%%%%%%%%%%%%%%%%%%%%%%%%%%%%%%%
% Required packages:

\RequirePackage{lineno}
\RequirePackage{color}

% Load natbib unless told otherwise
\if@natbib
  \RequirePackage[square,numbers]{natbib}
  \bibliographystyle{corlabbrvnat}
\fi

% set page geometry
\RequirePackage{hyperref}       % hyperlinks
\RequirePackage[verbose=true,letterpaper]{geometry}
\AtBeginDocument{
  \newgeometry{
    textheight=9in,
    textwidth=5.5in,
    top=1in,
    headheight=12pt,
    headsep=25pt,
    footskip=30pt
  }
  \@ifpackageloaded{fullpage}
    {\PackageWarning{corl_2026}{fullpage package not allowed! Overwriting formatting.}}
    {}
}


\ifdefined\nohyperref\else\ifdefined\hypersetup
  \definecolor{mydarkblue}{rgb}{0,0.08,0.45}
  \hypersetup{ %
    pdftitle={},
    pdfauthor={},
    pdfsubject={Proceedings of the \@conferenceordinal\/ Conference on Robot Learning (CoRL \@conferenceyear)},
    pdfkeywords={},
    pdfborder=0 0 0,
    pdfpagemode=UseNone,
    colorlinks=true,
    linkcolor=mydarkblue,
    citecolor=mydarkblue,
    filecolor=mydarkblue,
    urlcolor=mydarkblue,
    pdfview=FitH}

  \ifdefined\isaccepted \else
    \hypersetup{pdfauthor={Anonymous Submission}}
  \fi
\fi\fi


%%%%%%%%%%%%%%%%%%%%%%%%%%%%%%%%%%%%%%%%%%%%%%%%%%%%%%%%%%%%%%%%%%%%%%%%
% 

% fonts
\renewcommand{\rmdefault}{ptm}
\renewcommand{\sfdefault}{phv}

% Create acknowledgments -- only if the option 'final' is activated
\providecommand{\acknowledgments}{}
\renewcommand{\acknowledgments}[1]{%
\if@conferencefinal% 
\subsubsection*{Acknowledgments} #1
\fi
\if@preprinttype% 
\subsubsection*{Acknowledgments} #1
\fi
}


% handle tweaks for camera-ready copy vs. submission copy
\if@conferencefinal
  \newcommand{\@noticestring}{%
    \@conferenceordinal\/ Conference on Robot Learning
    (CoRL \@conferenceyear), \@conferencelocation.%
  }
\else
  \if@preprinttype
    \newcommand{\@noticestring}{%
      % Nothing here.
    }
  \else
    \newcommand{\@noticestring}{%
      Submitted to the \@conferenceordinal\/ Conference on Robot Learning (CoRL \@conferenceyear). Do not distribute.%
    }

    % line numbers for submission
    \linenumbers

    % fix incompatibilities between lineno and amsmath, if required, by
    % transparently wrapping linenomath environments around amsmath
    % environments
    \AtBeginDocument{%
     \@ifpackageloaded{amsmath}{%
        \newcommand*\patchAmsMathEnvironmentForLineno[1]{%
          \expandafter\let\csname old#1\expandafter\endcsname\csname #1\endcsname
          \expandafter\let\csname oldend#1\expandafter\endcsname\csname end#1\endcsname
          \renewenvironment{#1}%
            {\linenomath\csname old#1\endcsname}%
            {\csname oldend#1\endcsname\endlinenomath}%
        }%
        \newcommand*\patchBothAmsMathEnvironmentsForLineno[1]{%
          \patchAmsMathEnvironmentForLineno{#1}%
          \patchAmsMathEnvironmentForLineno{#1*}%
        }%
        \patchBothAmsMathEnvironmentsForLineno{equation}%
        \patchBothAmsMathEnvironmentsForLineno{align}%
        \patchBothAmsMathEnvironmentsForLineno{flalign}%
        \patchBothAmsMathEnvironmentsForLineno{alignat}%
        \patchBothAmsMathEnvironmentsForLineno{gather}%
        \patchBothAmsMathEnvironmentsForLineno{multline}%
      }{}
    }
  \fi
\fi



% The DOI will now automatically generate a URL link. Pretty cool!
% + Fix for hyperref and DOI (https://www.tug.org/pipermail/tex-live/2012-August/032161.html)
%-----------------------------
\makeatletter
\providecommand{\doi}[1]{%
  \begingroup
    \let\bibinfo\@secondoftwo
    \urlstyle{rm}%
    \href{http://dx.doi.org/#1}{%
      doi:\discretionary{}{}{}%
      \nolinkurl{#1}%
    }%
  \endgroup
}
% \makeatother
% %-----------------------------

\widowpenalty=10000
\clubpenalty=10000
\flushbottom
\sloppy

% font sizes with reduced leading
\renewcommand{\normalsize}{%
  \@setfontsize\normalsize\@xpt\@xipt
  \abovedisplayskip      7\p@ \@plus 2\p@ \@minus 5\p@
  \abovedisplayshortskip \z@ \@plus 3\p@
  \belowdisplayskip      \abovedisplayskip
  \belowdisplayshortskip 4\p@ \@plus 3\p@ \@minus 3\p@
}
\normalsize
\renewcommand{\small}{%
  \@setfontsize\small\@ixpt\@xpt
  \abovedisplayskip      6\p@ \@plus 1.5\p@ \@minus 4\p@
  \abovedisplayshortskip \z@  \@plus 2\p@
  \belowdisplayskip      \abovedisplayskip
  \belowdisplayshortskip 3\p@ \@plus 2\p@   \@minus 2\p@
}
\renewcommand{\footnotesize}{\@setfontsize\footnotesize\@ixpt\@xpt}
\renewcommand{\scriptsize}{\@setfontsize\scriptsize\@viipt\@viiipt}
\renewcommand{\tiny}{\@setfontsize\tiny\@vipt\@viipt}
\renewcommand{\large}{\@setfontsize\large\@xiipt{14}}
\renewcommand{\Large}{\@setfontsize\Large\@xivpt{16}}
\renewcommand{\LARGE}{\@setfontsize\LARGE\@xviipt{20}}
\renewcommand{\huge}{\@setfontsize\huge\@xxpt{23}}
\renewcommand{\Huge}{\@setfontsize\Huge\@xxvpt{28}}

% sections with less space
\providecommand{\section}{}
\renewcommand{\section}{%
  \@startsection{section}{1}{\z@}%
                {-2.0ex \@plus -0.5ex \@minus -0.2ex}%
                { 1.5ex \@plus  0.3ex \@minus  0.2ex}%
                {\large\bf\raggedright}%
}
\providecommand{\subsection}{}
\renewcommand{\subsection}{%
  \@startsection{subsection}{2}{\z@}%
                {-1.8ex \@plus -0.5ex \@minus -0.2ex}%
                { 0.8ex \@plus  0.2ex}%
                {\normalsize\bf\raggedright}%
}
\providecommand{\subsubsection}{}
\renewcommand{\subsubsection}{%
  \@startsection{subsubsection}{3}{\z@}%
                {-1.5ex \@plus -0.5ex \@minus -0.2ex}%
                { 0.5ex \@plus  0.2ex}%
                {\normalsize\bf\raggedright}%
}
\providecommand{\paragraph}{}
\renewcommand{\paragraph}{%
  \@startsection{paragraph}{4}{\z@}%
                {1.5ex \@plus 0.5ex \@minus 0.2ex}%
                {-1em}%
                {\normalsize\bf}%
}
\providecommand{\subparagraph}{}
\renewcommand{\subparagraph}{%
  \@startsection{subparagraph}{5}{\z@}%
                {1.5ex \@plus 0.5ex \@minus 0.2ex}%
                {-1em}%
                {\normalsize\bf}%
}
\providecommand{\subsubsubsection}{}
\renewcommand{\subsubsubsection}{%
  \vskip5pt{\noindent\normalsize\rm\raggedright}%
}

% float placement
\renewcommand{\topfraction      }{0.85}
\renewcommand{\bottomfraction   }{0.4}
\renewcommand{\textfraction     }{0.1}
\renewcommand{\floatpagefraction}{0.7}

\newlength{\@nipsabovecaptionskip}\setlength{\@nipsabovecaptionskip}{7\p@}
\newlength{\@nipsbelowcaptionskip}\setlength{\@nipsbelowcaptionskip}{\z@}

\setlength{\abovecaptionskip}{\@nipsabovecaptionskip}
\setlength{\belowcaptionskip}{\@nipsbelowcaptionskip}

% swap above/belowcaptionskip lengths for tables
\renewenvironment{table}
  {\setlength{\abovecaptionskip}{\@nipsbelowcaptionskip}%
   \setlength{\belowcaptionskip}{\@nipsabovecaptionskip}%
   \@float{table}}
  {\end@float}

% footnote formatting
\setlength{\footnotesep }{6.65\p@}
\setlength{\skip\footins}{9\p@ \@plus 4\p@ \@minus 2\p@}
\renewcommand{\footnoterule}{\kern-3\p@ \hrule width 12pc \kern 2.6\p@}
\setcounter{footnote}{0}

% paragraph formatting
\setlength{\parindent}{\z@}
\setlength{\parskip  }{5.5\p@}

% list formatting
\setlength{\topsep       }{4\p@ \@plus 1\p@   \@minus 2\p@}
\setlength{\partopsep    }{1\p@ \@plus 0.5\p@ \@minus 0.5\p@}
\setlength{\itemsep      }{2\p@ \@plus 1\p@   \@minus 0.5\p@}
\setlength{\parsep       }{2\p@ \@plus 1\p@   \@minus 0.5\p@}
\setlength{\leftmargin   }{3pc}
\setlength{\leftmargini  }{\leftmargin}
\setlength{\leftmarginii }{2em}
\setlength{\leftmarginiii}{1.5em}
\setlength{\leftmarginiv }{1.0em}
\setlength{\leftmarginv  }{0.5em}
\def\@listi  {\leftmargin\leftmargini}
\def\@listii {\leftmargin\leftmarginii
              \labelwidth\leftmarginii
              \advance\labelwidth-\labelsep
              \topsep  2\p@ \@plus 1\p@    \@minus 0.5\p@
              \parsep  1\p@ \@plus 0.5\p@ \@minus 0.5\p@
              \itemsep \parsep}
\def\@listiii{\leftmargin\leftmarginiii
              \labelwidth\leftmarginiii
              \advance\labelwidth-\labelsep
              \topsep    1\p@ \@plus 0.5\p@ \@minus 0.5\p@
              \parsep    \z@
              \partopsep 0.5\p@ \@plus 0\p@ \@minus 0.5\p@
              \itemsep \topsep}
\def\@listiv {\leftmargin\leftmarginiv
              \labelwidth\leftmarginiv
              \advance\labelwidth-\labelsep}
\def\@listv  {\leftmargin\leftmarginv
              \labelwidth\leftmarginv
              \advance\labelwidth-\labelsep}
\def\@listvi {\leftmargin\leftmarginvi
              \labelwidth\leftmarginvi
              \advance\labelwidth-\labelsep}

% create title
\providecommand{\maketitle}{}
\renewcommand{\maketitle}{%
  \par
  \begingroup
    \renewcommand{\thefootnote}{\fnsymbol{footnote}}
    % for perfect author name centering
    \renewcommand{\@makefnmark}{\hbox to \z@{$^{\@thefnmark}$\hss}}
    % The footnote-mark was overlapping the footnote-text,
    % added the following to fix this problem               (MK)
    \long\def\@makefntext##1{%
      \parindent 1em\noindent
      \hbox to 1.8em{\hss $\m@th ^{\@thefnmark}$}##1
    }
    \thispagestyle{empty}
    \@maketitle
    \@thanks
    \@notice
  \endgroup
  \let\maketitle\relax
  \let\thanks\relax
}

% rules for title box at top of first page
\newcommand{\@toptitlebar}{
  \hrule height 4\p@
  \vskip 0.25in
  \vskip -\parskip%
}
\newcommand{\@bottomtitlebar}{
  \vskip 0.29in
  \vskip -\parskip
  \hrule height 1\p@
  \vskip 0.09in%
}

%% keywords as first class citizens
\def\keywords#1{%
%  \ifdefined\isaccepted \else
%    \par {\bf Keywords:} #1%
%  \fi
%  \ifdefined\nohyperref\else\ifdefined\hypersetup
%    \hypersetup{pdfkeywords={#1}}
%  \fi\fi
 \ifdefined\isaccepted \else
    \begin{quote}
    \textbf{Keywords:} #1%
    \end{quote}
 \fi
  \ifdefined\nohyperref\else\ifdefined\hypersetup
    \hypersetup{pdfkeywords={#1}}
  \fi\fi
}


% create title (includes both anonymized and non-anonymized versions)
\providecommand{\@maketitle}{}
\renewcommand{\@maketitle}{%
  \vbox{%
    \hsize\textwidth
    \linewidth\hsize
    \vskip 0.1in
%     \@toptitlebar
    \centering
    {\LARGE\bf \@title\par}
%     \@bottomtitlebar
    \if@conferencefinal
      \def\And{%
        \end{tabular}\hfil\linebreak[0]\hfil%
        \begin{tabular}[t]{c}\bf\rule{\z@}{24\p@}\ignorespaces%
      }
      \def\AND{%
        \end{tabular}\hfil\linebreak[4]\hfil%
        \begin{tabular}[t]{c}\bf\rule{\z@}{24\p@}\ignorespaces%
      }
      \begin{tabular}[t]{c}\bf\rule{\z@}{24\p@}\@author\end{tabular}%
    \else
      \if@preprinttype
        \def\And{%
          \end{tabular}\hfil\linebreak[0]\hfil%
          \begin{tabular}[t]{c}\bf\rule{\z@}{24\p@}\ignorespaces%
        }
        \def\AND{%
          \end{tabular}\hfil\linebreak[4]\hfil%
          \begin{tabular}[t]{c}\bf\rule{\z@}{24\p@}\ignorespaces%
        }
        \begin{tabular}[t]{c}\bf\rule{\z@}{24\p@}\@author\end{tabular}%
      \else
        \begin{tabular}[t]{c}\bf\rule{\z@}{24\p@}
          Anonymous Author(s) \\
          Affiliation \\
          Address \\
          \texttt{email} \\
        \end{tabular}%
      \fi
    \fi
    \vskip 0.3in \@minus 0.1in
  }
}

% add conference notice to bottom of first page
\newcommand{\ftype@noticebox}{8}
\newcommand{\@notice}{%
  % give a bit of extra room back to authors on first page
  \enlargethispage{2\baselineskip}%
  \@float{noticebox}[b]%
    \footnotesize\@noticestring%
  \end@float%
}

% abstract styling
\renewenvironment{abstract}%
{%
%   \vskip 0.075in%
%   \centerline%
%   {\large\bf Abstract}%
%   \vspace{0.5ex}%
  \begin{quote}%
  \textbf{Abstract:}%
}
{
  \par%
  \vskip 1ex%
%   \ifdefined\keywords
%   \textbf{Keywords:}%
%   \@keywords%
%   \else
%   \fi
  \end{quote}%
}

\endinput

\end{filecontents*}
\documentclass{article}
\usepackage[T1]{fontenc}
\usepackage{corl_2026}
\usepackage{amsmath,amssymb,booktabs}
\hypersetup{pdftitle={Auditing Native Grasp Transfer Before Learning Failure Prediction},pdfauthor={},pdfsubject={Anonymous workshop submission}}
\title{Auditing Native Grasp Transfer\\Before Learning Failure Prediction}
\author{Anonymous Authors}
\begin{document}
\maketitle
\begin{abstract}
Reliable grasp-failure prediction requires trustworthy transfer, execution, and supervision. We present a layered audit of HUG-generated grasps executed with the native MANO hand in MANOSim. The audit separates source-coordinate fidelity, compiled object geometry, component dynamics, and task outcomes. Twelve component fixtures and six timestep comparisons pass. Across a predeclared 30-object geometry-by-size cohort, compiled object vertices reproduce the authored assets within 15.3 nanometres. We identify an unsupported validation assumption that confused depth encoding precision with stereo and reconstruction accuracy, and explicitly restrict subsequent evaluation to the authored simulation assets. Of 120 primary task slots, 13 complete as failures and 107 are quarantined; all 15 matched timestep slots are also quarantined. Eleven eligible pre-lift histories remain, all failures, so discriminative predictor training is not supported. These preliminary findings show that accurate coordinate transfer does not ensure valid execution or trainable outcome diversity. Next steps address passive object settling, standoff approach, independent scene calibration, and grouped evaluation after collecting both outcome classes.
\end{abstract}
\keywords{Dexterous grasping, simulation validation, grasp transfer, failure prediction}

\section{Introduction}
A grasp-failure predictor can learn pipeline artifacts when labels conflate mapping errors, unstable initialization, controller limitations and completed task failures. More simulated trials do not resolve this ambiguity. We audit Human Universal Grasping (HUG)~\citep{hug}, which generates MANO poses from RGB-D, using its native hand in MANOSim~\citep{manosim} to avoid cross-anatomy retargeting.

We distinguish \emph{coordinate fidelity}, \emph{numerical execution} and \emph{physical scene fidelity}. Source/compiled geometry checks establish numerical realization; dynamics test actuation/contact behavior; real-scene agreement needs sensor/reconstruction calibration. Our contribution is a versioned validation and supervision audit, not a forecasting architecture. The resulting cohort has eleven eligible histories but no successes, so we report no classifier score.

\section{Source transfer and execution contract}
\paragraph{Native realization.}
We rebuild the right MANOSim hand at HUG's exported shape, retaining authored masses, inertias, actuator settings and capsule/palm colliders. Its 67 positions, 51 velocities and 51 controls become 74 positions and 57 velocities with a free object. Capsules approximate MANO skin; they do not establish full-surface contact equivalence~\citep{mano,manosim}.

We retain exported float32 MANO axis angles/translation and the original float64 camera-to-world transform. Native commands account for authored rest rotations and project principal rotation vectors onto actuator axes. Static checks cover anchors, all sixteen bone orientations, actuator lengths and zero-error servo forces. Intrinsics must match the exact upstream float32 cast; correcting an initial float64-equality check changes no camera transform.

Objects retain mesh scale, mass, inertia and COM placement. Restoring compiled vertices through actual compiler recentering/principal frames must reproduce literal float32 OBJ vertex coverage bidirectionally within the pre-existing $5\,\mu$m bound. This checks vertex frames, not independently replayed topology, hull construction or physical reconstruction accuracy.

\paragraph{Clocks and outcome separation.}
Commands/histories are recorded every 10\,ms; primary physics uses 2\,ms steps. MuJoCo~\citep{mujoco} post-step contact/body diagnostics describe the preceding interval; position/velocity samples are endpoints. Contact/object diagnostics supply supervision only. Frozen protocols bind inputs, sources, runtime and output bytes. Independent readers replay saved geometry, clocks, labels and contact arithmetic without new generation/physics. Mapping refusal and operating-envelope/runtime quarantine remain null, separate from completed task failure.

\paragraph{Variety and partitions.}
HUG-Bench supplies five geometry categories $\times$ three authored sizes $\times$ two test objects~\citep{hug}. Two fixed captures and seeds 42/43 yield 120 primary proposals; fifteen 1\,ms re-evaluations cover one fixed object/capture/seed per cell. HUG test objects are not automatically unseen to our predictor. UniDexGrasp~\citep{unidex} motivates separating pose diversity from object/category generalization: all captures, seeds and scenarios of each object stay together in two reciprocal object and five category partitions. These support readiness analysis; the trainer separately implements five whole-category folds with a validation category. No fits occurred. Dex-Net~\citep{dexnet} motivates later physics/sensing variation, not replacing executed outcomes with analytic labels. Noise and nominal parameters remain fixed here.

\section{Validation findings}
\paragraph{Unsupported depth admission.}
An initial preflight backprojected 32 deterministic pixels per capture from a two-pixel-eroded mask. Each point had to lie within one depth bin plus a half-pixel diagonal footprint of the authored visual mesh. All 60 captures failed: only 837/1,920 points met their allowances. Intrinsics correction left these refusals unchanged; both preflight versions integrated zero task steps.

Encoding precision does not bound stereo, segmentation, SLAM or reconstruction error. HUG reprojects/crops estimated stereo depth, fills holes, and pose/scale-edits, watertights and decimates reconstructed meshes~\citep{hug,aria2mesh}. The intrinsics cast moves sampled points by at most 33.1\,nm, insufficient to explain millimetre residuals.

Exploratory fixed-transform checks give median silhouette recall 98.45\%, median IoU 81.84\%, and median per-view absolute ray-depth residual 8.39\,mm (7,054 hits/7,265 interior rays). No pose/scale was fitted. These are coarse correspondence diagnostics, not calibrated surface accuracy. A separate pilot therefore studies outcomes \emph{conditional on authored simulation assets}; physical RGBD correspondence remains unqualified and failed preflights are preserved.

\paragraph{Bounded dynamics.}
Two development poses $\times$ three fixtures $\times$ two timesteps give twelve tests: gravity motion/hold, simultaneous index/middle plane contacts, and sphere/table approach checks covering all sixteen components. Hand-body gravity is compensated; object gravity, moving Coriolis terms, targets, gains and inertias remain unchanged. The authored uncompensated baseline shows 13.08\,mm sag; compensated conditions meet fixed tracking gates (Table~\ref{tab:validation}). These different motion/hold conditions are not a hardware comparison.

Earlier fixtures missed contacts after runtime plane motion and exceeded a 2\,mm moving-window goal limit. New versions author planes before compilation and add approach settling, retaining limits and failures. Collision-search arrays remain unchanged. All twelve corrected fixtures and six timestep pairs pass; both finger reactions are positive throughout assessed tails. Actual clearance is sampled per step, not continuously certified for arbitrary servo motion.

\begin{table}[t]
\centering
\caption{Finite validation evidence. Distances concern numerical realization or tracking; they do not measure real-world accuracy.}
\label{tab:validation}
\begin{tabular}{lr}\toprule
Check & Observed maximum / result\\\midrule
Source anchor discrepancy, object pilot & $0.3291\,\mu$m\\
Source rotation-component discrepancy & $3.7942\times10^{-7}$\\
Compiled object vertex discrepancy, 432 instances & $15.3$\,nm\\
Gravity settled wrist error & $1.26\,\mu$m\\
Approach settled wrist error & $3.65\,\mu$m\\
Two-finger fixture penetration & $0.3271$\,mm\\
Component fixtures / timestep pairs & $12/12$; $6/6$ pass\\\bottomrule
\end{tabular}
\end{table}

\section{Object-task pilot and learning readiness}
\paragraph{Task and guards.}
The scene contains one target and a table, omitting source-image clutter. Closing occupies 1.0--1.5\,s, the wrist lifts 0.15\,m during 1.5--2.5\,s, and holding ends at 5.5\,s. Success requires COM rise $\geq0.10$\,m at lift end and $\geq0.08$\,m throughout holding, wrist-relative COM shift $\leq0.02$\,m, and no disallowed hand/table contact or hold-time object/table support. Contact/support events use aggregate normal force strictly above 1\,mN at every physics tick.

The earliest event determines failure. Hand--object contact before closing or hand--table contact at any phase is labeled approach collision, retaining actual phase. Missing acquisition contact, failed acquisition, drop and hold-relative-shift events are distinct; the last is a slip proxy, not tactile onset. Quarantine guards cover no-contact reset drift above 1\,mm/2 degrees, preclose tracking above 2\,mm/2 degrees, principal-chart violations and penetration above 5\,mm. Envelope violations do not themselves prove invalid engine arithmetic.

\paragraph{Outcomes.}
All 135 source/scene gates pass. The pilot uses 45 model/data pairs, 283 static forwards, 270 resets, 89,915 integration steps and 437,953 force reads. Independent readers verify 1,362 creation-time artifacts and 2,196 input bindings. Across 432 mesh instances/45 models, vertex discrepancy is at most 15.3\,nm. At 218 contact witnesses, 57-coordinate Jacobian and virtual-work errors are at most $4.16\times10^{-17}$ and $6.66\times10^{-16}$\,W: local numerical checks, not full-stream or hardware force calibration.

All thirteen completed primary attempts fail (Table~\ref{tab:outcomes}). All fifteen matched re-evaluations quarantine, so task-level timestep convergence is \emph{not established}; the six passing pairs concern separate fixtures. Nulls remain unlabeled, and completed-subset outcomes do not estimate failure rate over all proposals.

\begin{table}[t]
\centering
\caption{Predeclared task-slot accounting. Quarantines have null physical outcomes. No completed success is observed.}
\label{tab:outcomes}
\begin{tabular}{lrr}\toprule
Status / first cause & Primary & Matched 1\,ms\\\midrule
Completed: failed acquisition & 8 & 0\\
Completed: approach collision & 2 & 0\\
Completed: no contact & 2 & 0\\
Completed: hold-relative-shift slip & 1 & 0\\\midrule
Null: reset support stability & 49 & 6\\
Null: penetration envelope & 36 & 4\\
Null: preclose tracking & 22 & 5\\\midrule
Total scheduled & 120 & 15\\\bottomrule
\end{tabular}
\end{table}

\paragraph{Learning readiness.}
Twenty closing endpoints end at 1.49\,s, one control interval before lift. Each has 67 positions, 51 velocities and 51 effective commands: 169 channels. Future/object/contact state, labels, identities and timestamps are excluded. Planned models are prevalence, regularized logistic regression, a small MLP and a small GRU, with training-only normalization and validation-only selection.

Excluding two already-observed approach failures leaves eleven eligible primary prefixes across six objects/two categories, all future failures. Every catalog and trainer fold is empty or single-class in at least one role. Scientific fitting is withheld; no discriminative accuracy, calibration, useful warning horizon or scaling benefit is established. Manufactured controls cannot supply missing outcome diversity.

\section{Limitations and next steps}
This finite coordinate/asset audit does not qualify physical transfer. Capsules approximate skin; compliance and the 5\,mm envelope are not calibrated tissue models. Clutter is omitted; physical RGBD correspondence remains uncalibrated. Views/seeds are correlated, and 30 objects do not establish open-world generalization. The inspected depth/silhouette diagnostics are exploratory, not untouched validation. No physical experiment or recovery controller is evaluated.

Next, settle objects with the hand clear, retain measured state, and re-express wrist placement by the rigid object delta while preserving local finger pose. Approach should stop at a standoff before coordinated translation/closure. These drafts are unvalidated and unexecuted. New versions must preserve explicit limits, null accounting and paired checks.

Calibrate scene correspondence on disjoint validation scenes, freeze visibility/outlier handling, and test held-out views. Collect both outcome classes and adequate object coverage before grouped learning; vary mass, effective friction, approach and sensing without leaking held-out objects. Faithful transfer is necessary but insufficient for trustworthy prediction data.

\paragraph{AI assistance.} AI tools assisted code, analysis and drafting. The authors are responsible for checking the evidence, interpretation, citations and submission.

\clearpage
\begin{thebibliography}{9}
\bibitem{hug} K. Y. Wu et al. Human Universal Grasping. arXiv:2606.17054, 2026. \url{https://arxiv.org/abs/2606.17054}.
\bibitem{manosim} K. Y. Wu. MANOSim: MANO hand models in MJCF format for MuJoCo simulation. Software, 2026. \url{https://github.com/KevinyWu/manosim}.
\bibitem{mano} J. Romero, D. Tzionas, and M. J. Black. Embodied Hands: Modeling and Capturing Hands and Bodies Together. \emph{ACM Transactions on Graphics}, 36(6), 2017. \url{https://doi.org/10.1145/3130800.3130883}.
\bibitem{unidex} Y. Xu et al. UniDexGrasp: Universal Robotic Dexterous Grasping via Learning Diverse Proposal Generation and Goal-Conditioned Policy. \emph{CVPR}, 2023. \url{https://arxiv.org/abs/2303.00938}.
\bibitem{dexnet} J. Mahler et al. Dex-Net 2.0: Deep Learning to Plan Robust Grasps with Synthetic Point Clouds and Analytic Grasp Metrics. \emph{Robotics: Science and Systems}, 2017. \url{https://arxiv.org/abs/1703.09312}.
\bibitem{aria2mesh} K. Y. Wu. aria2mesh: Metrically accurate 3D object reconstruction and simulation-ready assets from Aria Gen 2 recordings. Software, 2026. \url{https://github.com/KevinyWu/aria2mesh}.
\bibitem{mujoco} E. Todorov, T. Erez, and Y. Tassa. MuJoCo: A physics engine for model-based control. \emph{IROS}, 2012. \url{https://doi.org/10.1109/IROS.2012.6386109}.
\end{thebibliography}
\end{document}
